% =====================================================================
%  TYPE: Implant MRI-compatibility certificate
%  For the imaging centre, before an MRI in a patient with implants.
%  Based on 'mri certification-gr'.
%
%  Copy, rename to "YYYYMMDD - surname - topic.tex", fill in and
%  compile with XeLaTeX. Toggle alternatives by moving the "%".
%  Text in [brackets] is a placeholder: replace or delete it.
% =====================================================================
\documentclass[10pt,greek]{article}
\usepackage[useregional]{datetime2}

%%%%%%%%%%%%%%%%%%%%% Language
\providecommand\DocLang{gr}

%%%%%%%%%%%%%%%%%%%%% Clinic (letterhead logo, footer address)
\providecommand\Clinic{Peristeri}
% \providecommand\Clinic{Mandra}
% \providecommand\Clinic{Bioclinic}
% \providecommand\Clinic{Psychikou}
% \providecommand\Clinic{Metropolitan}
% \providecommand\Clinic{Dunant}
% \providecommand\Clinic{Euroclinic}

%%%%%%%%%%%%%%%%%%%%% Document code for the QR: paste the code issued
%%%%%%%%%%%%%%%%%%%%% by doc.erns.eu (empty = no QR code)
\providecommand\DocCode{}

%%%%%%%%%%%%%%%%%%%%% Patient
\providecommand\PatientPrefix{Ο}
% \providecommand\PatientPrefix{Η κα}
\providecommand\PatientName{Παπαδοπούλου Μαρία}
\providecommand\PatientAge{45}
% \providecommand\PatientAge{}

%%%%%%%%%%%%%%%%%%%%% Date
\providecommand\DocDate{\DTMdate{2026-01-15}}
% \providecommand\DocDate{\DTMdate{2026-05-04}}

%%%%%%%%%%%%%%%%%%%%% Title
\providecommand\DocTitle{Ιατρική Βεβαίωση}

%%%%%%%%%%%%%%%%%%%%% Closing (empty = no signature block)
\providecommand\DocClosing{Ο ιατρός}
% \providecommand\DocClosing{}

%%%%%%%%%%%%%%%%%%%%% Stamp and signature image
\providecommand\DocStamp{no}
% \providecommand\DocStamp{yes}

\input{evrtemplate}

\begin{document}

\begin{PatientInfo}
    \Info{Ημ. Χειρουργείου}{10-01-2026}
\end{PatientInfo}

Η ασθενής υποβλήθηκε στις 10/01/2026 σε διαδερμική σπονδυλοδεσία ΟΜΣΣ με χρήση νευροπλοηγού.

Βεβαιώνεται ότι τα υλικά που χρησιμοποιήθηκαν είναι από παραμαγνητικό υλικό, και συγκεκριμένα τιτάνιο / PEEK, που είναι απόλυτα συμβατά με Μαγνητικό Τομογράφο. Επομένως η ασθενής μπορεί να υποβληθεί σε μαγνητική τομογραφία με ασφάλεια.

% Purpose (for a female patient: "της ασθενούς", "της")
% Το παρόν χορηγείται για χρήση στην ασφαλιστική εταιρεία του ασθενούς.
% Το παρόν χορηγείται για χρήση από τον ασφαλιστικό φορέα του ασθενούς.
Το παρόν χορηγείται για την εργασία της ασθενούς.
% Το παρόν χορηγείται μετά από αίτηση του ασθενούς, για την υπηρεσία του.
% Το παρόν χορηγείται για χρήση στο διαγνωστικό κέντρο.
% Το παρόν χορηγείται για κάθε νόμιμη χρήση.
% Η παρούσα βεβαίωση δίνεται κατόπιν επιθυμίας του ασθενούς.

\end{document}
